% RSS Stage 1: keep the supplied conference layout. Compile from this directory.
\documentclass[conference]{IEEEtran}
\usepackage{times}
\usepackage[numbers]{natbib}
\usepackage{amsmath,booktabs}
\usepackage[bookmarks=true]{hyperref}
\hypersetup{pdfauthor={Anonymous},pdftitle={When Should a Robot Update Its Digital Twin?}}
\begin{document}
\title{When Should a Robot Update Its Digital Twin?\\Choosing Learning Operations from Deployment Feedback}
\author{Author names omitted for anonymous review}
\maketitle
\begin{abstract}
A missed placement can indicate an incorrect estimate of tray height or a layout absent from the demonstrations. The first suggests changing the simulator; the second suggests changing the data generated in it. We study this choice in a controlled MuJoCo packing task, with separate training and deployment scenes. In 72 fixed-recipe runs, replacing one round of dataset aggregation with height calibration increased success by 33.3 percentage points under height error but reduced it by 21.7 points under limited layout coverage. In a separate 288-run comparison, a rule using deployment residuals and failure counts improved mean success over a calibrate-once schedule by 13.5 points (paired bootstrap interval: 8.6 to 18.6). The evidence used to accept a repair also mattered: a check based on the current policy's success inside the twin rejected seven of nine geometric updates, whereas a held-out height-prediction check rejected none and preserved deployment performance. These results distinguish two decisions: which learning operation to try, and how to test whether it repaired the model. They concern controlled simulation, not yet physical transfer.
\end{abstract}
\IEEEpeerreviewmaketitle
\section{Introduction}
Consider a robot that repeatedly releases an object above its slot. If the simulator places the tray too low, more demonstrations generated on that tray may teach the wrong release height more accurately. If the tray is correctly modeled but training covered only the left half of the workspace, measuring its height again supplies no examples on the right. Both failures can produce the same binary outcome. They differ in what the next deployment should buy: a better model of the scene or better coverage of the policy's inputs.

This distinction matters under a short improvement budget. Dataset aggregation (DAgger) labels states visited by a learner \cite{ross2011dagger}; a calibration changes the geometry used by the planner and the policy's state encoder. Spending a round on one leaves fewer opportunities for the other. Conversely, repairing the twin may change the value of demonstrations generated \emph{after} the repair. The relevant decision is therefore not just whether simulation helps, but which part of the learning pipeline to change next.

We investigate three questions. Does the preferred fixed recipe change when the failure source changes but the task does not? Can deployment observations guide a choice among those recipes' constituent operations? Finally, can the policy itself verify a proposed correction to the simulator in which it was trained? Our comparisons distinguish final task success from agreement with a prescribed first action, and both from the number of labels consumed. In particular, stopping can be the intended response to a healthy twin even if another training round would increase success.

Our contribution is a controlled comparison rather than a new low-level controller. We vary the relation between the training twin, the demonstration distribution and the deployment scene; compare recipes and selectors on matched task sets; and isolate a failure mode of policy-based model validation. We also specify an order intervention that could refute the need for adaptive routing: if one fixed sequence works across causes at the same cost, it is a credible alternative.
\clearpage
\section{Related Work}
\textbf{What to update.} RialTo reconstructs a scene, transfers demonstrations into the twin, and uses simulation reinforcement learning to improve a real-world imitation policy \cite{rialto}. MimicGen adapts demonstrations to new configurations \cite{mandlekar2023mimicgen}; ARSTAG uses agents to generate task-specific robot data \cite{arstag}. Agentic RSR already uses execution feedback in scene construction and policy development \cite{agenticrsr}. Our distinction is not the use of feedback itself. We hold the packing task fixed and compare spending a post-deployment round on the scene model against spending it on demonstrations or policy training.

\textbf{What to label.} DAgger changes the state distribution on which an expert supplies actions; HG-DAgger lets a human decide when to intervene \cite{ross2011dagger,kelly2019hgdagger}. LoopWAM is especially close: it collects deployment trajectories, annotates quality and retrains with expert and historical rollout data. On RoboChallenge Table30-V2, success rose from 23.67\% to 43.33\% over three rounds \cite{loopwam}. That result establishes the value of learning from mixed-quality physical rollouts. It does not compare collecting another batch with spending that round correcting geometry. Our deployment-label arm addresses that comparison in simulation; it is not a reproduction of LoopWAM.

\textbf{What to measure.} Domain randomization varies the training distribution \cite{tobin2017dr}. SimOpt adapts that distribution using real experience \cite{chebotar2019asid}, and ASID selects informative interactions for system identification \cite{asid}. These approaches motivate alternatives to assuming a known model. REFLECT summarizes executions to explain failures and guide replanning \cite{reflect}; AHA uses failure explanations to refine rewards, plans and verification \cite{aha}; CycleVLA backtracks at the subtask level \cite{cyclevla}. Rather than evaluate the explanation or the next motion, we measure the return on the next round of model or policy updates.

\section{Decision and Evaluation Protocol}
Let $W$ be a deployment scene and $T$ its reconstruction. A policy $\pi$ acts using features computed with $T$; success is scored in $W$. Both scenes use MuJoCo \cite{mujoco}. A bank is a set of task instances, not a robot. The task is to pack four objects in sequence using a hierarchical controller and a 55-dimensional state representation. The selector receives no cause identifier. Its inputs include observed support-height residuals, the distribution of failed subgoals, recent success and the history of updates. The injected cause is available only for analysis and to a reference selector.

At round $t$, a selector chooses an operation $o_t$ from this information. We compare final complete-task success rates (SR) under limits on feedback episodes, training transitions and rounds, recording actual expenditure rather than equating a label with a deployment. Standard initialization uses 80,000 labeled transitions. Each of three rounds provides 24 feedback episodes; a data operation adds up to 40,000 labels and retraining uses 8,000 optimization steps. Calibration uses observations but adds no training labels. Stopping saves subsequent rounds. The open-loop control spends its 200,000-label allowance at initialization. Feedback episodes are not themselves 40,000 expert labels, and deployment-side expert data would incur a different physical acquisition cost.

The geometric condition sets believed support height to 0.46 of the true height (9.2 versus 20 mm). Under limited coverage, training layouts occupy half the workspace and deployment uses the full layout distribution. The recovery condition omits recovery demonstrations and tests under perturbation; the healthy condition introduces neither a reconstruction error nor a training deficit. Healthy thus means no injected defect, not guaranteed success. Only recovery uses perturbed evaluation in the four-condition comparison. Differences between methods within a condition have a cleaner interpretation than differences between conditions.

\clearpage
\section{Which Repairs Help?}
\begin{table}[t]
\centering\caption{Fixed recipes: final deployment SR and mean additional labels (thousands, excluding initialization). Each SR averages three layout banks at one training seed, with 20 evaluation episodes per cell. The true-geometry arm is privileged.}\label{tab:recipe}
\begin{tabular}{@{}lrrrr@{}}\toprule
Recipe & Height & Coverage & Healthy & Labels\\\midrule
Open loop (200k initial) & .367 & .083 & .917 & 0\\
Fixed twin + DAgger & .417 & .933 & .950 & 102\\
Randomization + DAgger & .450 & .383 & 1.000 & 111\\
Deployment-side DAgger & .533 & .900 & .983 & 111\\
Mixed-source training & .283 & .867 & .983 & 120\\
Calibrate once + DAgger & .750 & .717 & .967 & 80\\
Alternate calibration/DAgger & .833 & .300 & .967 & 40\\
True geometry + DAgger & .950 & .933 & .950 & 93\\\bottomrule
\end{tabular}
\end{table}
Across 72 runs (Table~\ref{tab:recipe}), replacing the first DAgger round with calibration improved geometric success from .417 to .750. The paired gain was 33.3 percentage points (pp; 95\% bootstrap interval $[25.0,45.0]$). The same substitution reduced coverage success from .933 to .717 ($-21.7$ pp, $[-30.0,-5.0]$). Alternating calibration and DAgger amplified the contrast: $+41.7$ pp under height error, $-63.3$ pp under limited coverage. The intervals resample three matched banks, not independent training seeds. The evidence is a reversal among these recipes, not a theorem that every fixed schedule must lose.

Why not simply label deployment trajectories? Expert labeling in the deployment simulator reached .900 under limited coverage but .533 under height error. Randomizing physics and aggregating data reached .383 and .450, respectively. These are implemented alternatives, not benchmarks against physical teleoperation or all forms of domain randomization. When geometry was correct, the fixed-twin learner could already acquire labels on deployment layouts. Calibration removed a training opportunity without expanding that support.

\subsection{Separating model error from policy error}
We hold a scripted expert fixed and vary only its geometric belief. Let $A_2$ be its success in $W$ with true geometry and $A$ its success in $W$ with twin geometry. Let $C$ and $D$ be the learned policy's success in $W$ and in $T$, respectively, using the same twin belief. Only $A_2-A$ is independent of the learned policy. $A-C$ measures a policy shortfall relative to that expert, not a calibrated estimate of missing knowledge.
\begin{table}[t]
\centering\caption{Nominal probe success. Ten episodes per condition and controller; bank labels reuse the probe task seeds, so we do not pool them as independent replications.}\label{tab:probe}
\begin{tabular}{@{}lrrrrrr@{}}\toprule
Condition & $A_2$ & $A$ & $C$ & $D$ & $A_2-A$ & $D-C$\\\midrule
Geometry & 1.00 & .00 & .30 & .80 & 1.00 & .50\\
Coverage & 1.00 & 1.00 & .10 & .10 & .00 & .00\\
Recovery & 1.00 & 1.00 & 1.00 & 1.00 & .00 & .00\\
Healthy & 1.00 & 1.00 & .70 & .70 & .00 & .00\\\bottomrule
\end{tabular}
\end{table}
Under height error the policy succeeded eight times in ten inside its twin but only three times in deployment. Conversely, low in-twin success under limited coverage (.10) coexisted with perfect expert success in deployment using that same geometry. Thus a low policy score need not indict the twin, and a high one need not vindicate it. Nominal probes also missed the recovery deficit: perturbations lowered policy SR from 1.00 to .00 there, compared with .70 to .40 in the healthy control. These are diagnostic contrasts, not independent generalization tests of a new selector.

A calibration ablation tested the practical consequence. A gate rejected a candidate if the incumbent policy's in-twin SR fell by more than five points. Across three seeds it rejected seven of nine geometric updates; mean nominal deployment SR was .450, versus .717 without the gate.

The alternative checks the observable that the repair aims to predict. An object's underside provides a height measurement when it rests on the tray or when its commanded descent stalls with no object beneath it. Each episode contributes its median measurement. We fit the repair on alternating feedback episodes. On the other episodes, let $E(h)=|J|^{-1}\sum_{j\in J}|z_j-h|$, where $z_j$ is the measured height for that fixture. With at least two check observations, we reject a proposed height $h'$ if $E(h')>E(h)+2$ mm; fixtures with insufficient observations do not veto the update. The gate rejected none of the nine geometric candidates and one update on an already correct twin; mean geometric SR remained .717. The score does not use the policy's success, although the policy still determines which contacts are observed. A bias shared by fitting and checking measurements can thus escape detection.

\clearpage
\section{Selecting Operations from Feedback}
After checking the budget, the rule calibrates if observed and believed support heights differ by more than 4 mm. Otherwise, failures under perturbation can trigger recovery demonstrations. If feedback SR is at least .85 and failures are diffuse with low perturbation pressure, the rule stops. If more than half the failures concentrate on one subgoal and fewer than 20,000 additional labels have been spent, it requests targeted regeneration. This last gate counts labels across the run, not labels on the failing subgoal: one 40,000-label update can therefore rule out targeted regeneration later. The rule also stops after two rounds without improvement; its default is a DAgger round. These decisions use deployment observations, \emph{not} the privileged expert contrast in Table~\ref{tab:probe}. A separate multi-label selector predicts geometry, coverage and recovery components from failure profiles; its training set excludes the three evaluation banks. Neither selector receives the injected cause at decision time.

We compare these choices with a fixed calibrate-once schedule, random and round-robin schedules, an upper-confidence-bound bandit (UCB) \cite{auer2002ucb}, a contextual linear bandit (LinUCB) \cite{linucb}, and a reference that knows the injected cause and plays a prescribed sequence. UCB explores operations based on observed improvement; LinUCB also conditions their estimated value on the failure profile. These are adaptations to the same operation menu, not published scores from the cited bandit papers.
\begin{table}[t]
\centering\caption{Selectors on four conditions. Entries are mean final SR over three banks and three training seeds (288 runs). First counts matches to the cause-informed reference's first operation, out of 36.}\label{tab:selectors}
\begin{tabular}{@{}lrrrrrr@{}}\toprule
Selector & Geom. & Cov. & Rec. & Healthy & Mean & First\\\midrule
Calibrate once & .772 & .600 & .833 & .911 & .779 & 9\\
Random & .300 & .939 & .867 & .972 & .769 & 0\\
Round robin & .894 & .872 & .822 & .939 & .882 & 9\\
UCB & .961 & .983 & .817 & .978 & .935 & 9\\
Rule & .900 & .928 & .900 & .928 & .914 & 21\\
LinUCB & .939 & .917 & .356 & .961 & .793 & 9\\
Multi-label & .828 & .928 & .828 & .883 & .867 & 28\\
Cause-informed & .939 & .972 & .939 & .844 & .924 & 36\\\bottomrule
\end{tabular}
\end{table}
The rule improved mean SR by 13.5 pp over calibrating once, with a paired 95\% percentile interval of $[8.6,18.6]$ pp (5,000 resamples of matched cell differences). The largest gain was on coverage: the rule could acquire missing layouts immediately. UCB had the highest observed mean. Its difference from the rule was $+2.1$ pp $[-1.1,5.1]$: it gained 5--6 points on geometry, coverage and healthy conditions, but lost 8.3 points on recovery. We do not treat the interval crossing zero as evidence of equivalence. The cause-informed schedule's corresponding difference was $+1.0$ pp $[-2.2,4.0]$; it is a reference for action agreement, not a maximum achievable SR.

\textbf{Success versus cost.} The multi-label selector matched more reference first actions than the rule (28 versus 21), yet its mean final SR was 4.7 points lower (paired interval $[-8.5,-1.0]$). On healthy tasks it used an average of 22,222 new labels, versus 44,444 for the rule and 80,000 for calibrate-once. The reference spent none: it stopped after observing feedback. To compare speed, we counted feedback episodes up to the first attainment of 90\% of the condition-specific reference mean. Multi-label averaged 29.3 episodes \emph{among the 27 of 36 runs that attained the threshold}; the rule averaged 35.6 among 33 of 36. A smaller conditional mean is not evidence that a selector reaches the threshold more reliably.

\textbf{More than one defect.} When height error and coverage deficit coexist, both a fixed schedule and a rule can start correctly by calibrating. Their second decisions can differ: the rule can now regenerate data on missing layouts. On a separate mixed-condition screen (four mixtures, two banks, one seed), UCB attained mean SR .863, the rule .731, calibrate-once .637 and the reference .900. Offline, a multi-label diagnosis recovered at least one true component for all 18 mixed profiles and all components for 72\%; single-label mixed-class recall was 6\%. These are different experiments. Detecting one component motivates repairing it and measuring the residual failure again; the offline result does not attribute UCB's return to a classifier it never used.

\clearpage
\section{Order, Verification and Transfer}
\textbf{When would a fixed sequence suffice?} Differences between whole recipes do not tell us whether one pair of operations works in a fixed order for both causes. We first examined recorded decisions. Under limited coverage, targeted regeneration increased SR by .822 on average immediately after calibration ($n=18$ decisions), by .175 after DAgger ($n=10$), and by .037 after targeted regeneration ($n=15$). Under height error the corresponding changes after calibration and DAgger were .254 and $-.111$. These groups differ in starting SR, remaining headroom and which selector chose the action. Subtracting their means does not identify an effect of order.

The intervention needed is a paired comparison of sequences starting from the same policy and belief. For example, hold the number of calibrations and data updates fixed and estimate final SR for calibrate--then--regenerate minus regenerate--then--calibrate on the same bank. Repeating one operation twice tests whether either operation alone suffices. If one ordering works well for both causes, the evidence would favor that fixed schedule over paying for diagnosis. We report no intervention estimate here.

\textbf{Does a language model help?} An API-backed selector received a structured deployment profile and returned an operation name; it did not control the arm. In a 32-run screen on single causes (four variants, four conditions, two banks, one seed), the full menu yielded mean SR .900, versus .850 with operator-supplied experience, .781 with a restricted menu and .806 with a cheaper model; withholding an operation changes what any planner can achieve. On mixed causes the comparison is not favourable to the language model: across three mixtures, two banks and one seed, the rule reached .733, the model .667 (paired difference $-6.7$ pp, $[-25.0,10.0]$ over six matched cells) and a variant carrying per-run memory .525; the cause-informed reference reached .892. Billed charges were \$0.0145 per cell for three calls. The failure is specific: in the three-cause mixture the third-round action that helps is adding recovery data (.650 $\rightarrow$ .800 with the reference), while the model repeated the targeted regeneration it had already tried (.500 $\rightarrow$ .200). With nominal evaluation, however, perturbation pressure is below the rule's .5 threshold, so neither selector can observe that a recovery deficit remains. We therefore read this as a limit of observability rather than of planning, and we do not claim that the language model addresses causes the rule omits.

\textbf{Limits of the test bed.} The known defects are interventions in one task family, not a taxonomy of natural robot failures. A scripted expert with true geometry is a privileged diagnostic control unavailable on a physical system. Even the nonprivileged expert rollouts consume interactions. Before a probe-driven selector can claim budget efficiency, these interactions must be charged against its alternatives. The three-seed selector comparison is distinct from the single-seed recipe experiment, and neither is a nested evaluation across task families. Dynamics sweeps also restricted the question we could test: reducing object friction by a factor of 100 did not change measured success, while adding two steps of latency collapsed success. Neither produced a graded repairable condition here. This finding is specific to this grasp and transport task.

\textbf{External tests.} For the AgileX Piper, tag-based object poses would preserve the state interface. We plan to compare calibration and further imitation on paired placements, measure tray height independently, and record successes out of 20 on held-out placements as well as robot time. On LIBERO \cite{libero}, layout perturbations from LIBERO-PRO \cite{liberopro} would instead test a visual policy such as Diffusion Policy \cite{chi2023diffusion}. The critical control is to generate equal numbers of demonstrations with the original and the observation-updated layout priors. Neither hardware nor visual transfer is a reported result.
\begin{table}[t]
\centering\caption{External validation reserved for the completed study. Dashes denote unmeasured outcomes, not failures.}\label{tab:pending}
\begin{tabular}{@{}lll@{}}\toprule
Comparison & Piper: unseen layouts & Visual: shifted SR\\\midrule
No twin update & --- & ---\\
Randomization & --- & ---\\
Feedback-selected update & --- & ---\\\bottomrule
\end{tabular}
\end{table}
% Replace the empty panel only with paired, equal-budget intervention results.
\begin{figure}[t]
\centering\fbox{\rule{0pt}{2.5em}\rule{.87\columnwidth}{0pt}}
\caption{Reserved for the order intervention: paired difference in final SR between calibrate--regenerate and regenerate--calibrate, by bank and failure condition. No estimate is available yet.}
\end{figure}
\section{Conclusion}
In four-step simulated packing, reserving a round for calibration helped when geometry was wrong and hurt when demonstrations lacked coverage. Deployment feedback made this distinction useful even to a simple rule. The verification experiment identifies a separate risk: judging a candidate twin by the incumbent policy can reject a geometric correction precisely when the policy has adapted to the old belief. Comparing predicted and observed geometry tests the repair more directly. Whether these decisions justify their measurement cost on physical robots remains open.

\clearpage
\bibliographystyle{plainnat}
\bibliography{references}
\end{document}
